\newpage
\setChapterHeader{6}{结论与展望}
\section{结论与展望}
\zihao{-4}

\subsection{主要贡献}

首先，论文从传统数论证明出发，梳理了其在 Lean 4 中形式化表达的基本路径。Mordell 方程的证明并不只是把已有文字证明逐句翻译成代码，而是需要重新拆分其中的数学结构。本文重点分析了二次整数环中的元素分解、理想分解、类群下降以及参数化解判别等环节，并说明这些内容在形式化证明中如何被转化为定义、引理和可验证的推理步骤。

其次，论文在 Lean 4 与 Mathlib4 环境下完成了一个 Mordell 方程相关的形式化项目。整个项目约 2800 行代码，分布在 12 个源文件中，内容包括基础结构、代数性质、下降法、类群计算以及最终解定理等部分。通过这一实现，原本依赖自然语言叙述的数论论证被整理为能够由 Lean 检查的形式化代码。

最后，论文结合实际开发过程，对 AI 工具在 Lean 形式化中的作用进行了总结。实践表明，AI 可以在解释证明思路、拆分引理、生成代码草稿、查找相关定理和修改报错等方面提供帮助，但其结果并不能直接视为可靠证明。最终的正确性仍然依赖 Lean 编译器验证，也需要人工检查代码是否真正表达了预期的数学内容。基于这一过程，本文总结了 AI 辅助 Lean 形式化证明开发中的一些经验和限制。

\subsection{不足之处}

本文仍存在一些不足。

首先，形式化对象仍局限于 Mordell 方程的若干特殊负整数参数。虽然这些情形已经包含较丰富的代数数论结构，但尚不能代表一般 Mordell 方程或更广泛椭圆曲线整数点问题的全部复杂性。对于更一般的参数，类群计算、局部条件分析和下降过程可能需要更多理论工具与 Mathlib 支持。

其次，本文的形式化路线仍较大程度依赖已有数学证明和人工设计的证明结构。AI 能够辅助生成和修复代码，但尚未真正从自然语言问题出发独立完成完整证明发现、形式化建模和验证闭环。换言之，当前 AI 更适合作为高效的证明工程助手，而不是完全可靠的自动数学家。

再次，多模型实验受模型版本、提示词、上下文长度、可用工具和运行环境影响较大。由于 AI 系统更新速度很快，某一次实验结果只能反映特定时间和特定设置下的表现。后续若要得到更稳定的结论，需要建立更标准化的测试任务、记录格式和可复现实验环境。

最后，本文主要关注证明能否通过验证，对代码可维护性、证明长度、API 稳定性和自动化策略复用等工程指标讨论仍不充分。对于更大规模的 Lean 数学库建设而言，这些因素同样会直接影响形式化项目的长期价值。

\subsection{未来展望}

未来工作可以从以下几个方向继续展开。

第一，可以进一步扩展 Mordell 方程形式化的数学范围。除了本文处理的 \(d=-1,-2,-5,-6,-13\) 外，还可以尝试更多参数情形，或进一步研究一般 Mordell 方程、椭圆曲线整数点、Selmer 群计算和局部整体原则等更深层内容。这需要在 Lean 中补充更多代数数论、椭圆曲线和计算数论工具。

第二，可以继续向 Mathlib4 贡献相关基础工具。本文形式化过程中涉及二次整数环、范数、理想运算、Dedekind 整环和类群等内容，其中部分工具具有一定通用性。若能将这些基础引理整理为更规范、可复用的 Mathlib4 接口，将有助于后续更多代数数论命题的形式化。

第三，可以通过更清晰的工程结构和工作流改善 AI 生成代码的质量。未来的 AI 辅助形式化系统不应只生成单段证明代码，而应结合仓库结构、定理检索、编译反馈、证明状态和版本控制信息，形成更稳定的“生成—检查—修复—复核”流程，从而减少定理陈述偏移、无效代码和不可维护证明。

综上所述，本文通过 Mordell 方程整数解问题的 Lean 4 形式化案例，展示了人工智能辅助形式化数学的可行性和现实价值。AI 可以提高证明开发效率，帮助使用者跨过部分 Lean 工程门槛；而 Lean 的形式验证机制又为 AI 生成内容提供了严格的可靠性约束。随着 Mathlib4 的持续完善和 AI 证明智能体的发展，二者结合有望在未来数学研究、数学教育和复杂证明验证中发挥更加重要的作用。
